\documentclass[../main.tex]{subfiles}
\begin{document}
\section{Tier 2 --- bounded perturbations and products}
\label{sec:tier2}

One controlled step away from the exactly-solved Gaussian: multiply by a bounded factor, or
take products. Both are certified by a single template each (Holley--Stroock,
tensorization), and both degrade gracefully into the harder tiers when their hypothesis
($\osc(W)<\infty$, independence) breaks.

\begin{family}[Bounded perturbation of a Gaussian]
\label{fam:perturbed-gaussian-1d}
Let $\Pstar\propto e^{-U}$ on $\R$ with $U(x)=\tfrac{x^2}{2\sigma^2}+W(x)$ and
$\osc(W)\le A<\infty$. Then
\[
  \CP(\Pstar)\ \le\ \sigma^2\,e^{A}.
\]
\end{family}

\begin{proof}
The base measure $\nu=N(0,\sigma^2)$ has $\CP(\nu)=\sigma^2$ by
Family~\ref{fam:gaussian-covariance}, and $\Pstar=Z^{-1}e^{-W}\nu$. Holley--Stroock
(Theorem~\ref{thm:holley-stroock}) transfers the constant with the factor
$e^{\osc W}\le e^{A}$.
\end{proof}

\noindent A typical instance is the bounded oscillatory perturbation $W(x)=\eps\cos(\omega x)$,
with $\osc(W)\le2\eps$ and certified bound $\sigma^2 e^{2\eps}$. The estimate is sharp up to
the $e^{A}$ factor and applies verbatim to ``Gaussian prior $+$ bounded log-likelihood
perturbation'' posteriors. Once $W$ develops a deep well ($\osc(W)\to\infty$) the family
crosses into the separated regime of Section~\ref{sec:tier5} and the bound becomes vacuous.

\begin{family}[Products and conjugate independent priors]
\label{fam:products}
If $\Pstar=\bigotimes_{i=1}^D\mu_i$, then $\CP(\Pstar)=\max_i\CP(\mu_i)$ and
$\CLS(\Pstar)=\max_i\CLS(\mu_i)$.
\end{family}

\noindent Immediate from tensorization (Theorem~\ref{thm:tensorization}). Many independent
Bayesian priors/posteriors reduce this way to one-dimensional constants computable by the
Hardy criterion (Theorem~\ref{thm:hardy-1d}): e.g.\ independent
$\theta_j\sim N(0,\tau_j^2)$ gives $\CP=\CLS=\CTCI=\max_j\tau_j^2$. The catch is that
classical LSI and $T_2$ require Gaussian-type tails, so products of exponential- or
polynomial-tailed factors (Laplace, Student) inherit a finite Poincar\'e constant but
\emph{not} LSI/$T_2$ --- the tier boundary made precise in Section~\ref{sec:tier4}.
\end{document}
